\documentclass[10pt,a4paper]{article}
\usepackage[margin=1.5cm,top=1.3cm,bottom=1.5cm]{geometry}
\usepackage{times}
\usepackage{amsmath,amssymb}
\usepackage{booktabs,tabularx}
\usepackage[hidelinks]{hyperref}
\usepackage{enumitem}
\usepackage[most]{tcolorbox}
\usepackage{titlesec}
\usepackage{xcolor}
\setlist{itemsep=1pt,topsep=2pt,leftmargin=*}
% generated by paper/make_values.py -- do not edit
\newcommand{\aucG}{0.539}
\newcommand{\aucGHi}{0.557}
\newcommand{\aucGLo}{0.521}
\newcommand{\aucGSd}{0.002}
\newcommand{\aucLenC}{0.542}
\newcommand{\aucLenCHi}{0.561}
\newcommand{\aucLenCLo}{0.524}
\newcommand{\aucLenCSd}{0.002}
\newcommand{\aucPA}{0.564}
\newcommand{\aucPAHi}{0.582}
\newcommand{\aucPALo}{0.545}
\newcommand{\aucPASd}{0.006}
\newcommand{\aucPB}{0.589}
\newcommand{\aucPBHi}{0.607}
\newcommand{\aucPBLo}{0.571}
\newcommand{\aucPBSd}{0.005}
\newcommand{\aucPC}{0.600}
\newcommand{\aucPCHi}{0.619}
\newcommand{\aucPCLo}{0.582}
\newcommand{\aucPCSd}{0.004}
\newcommand{\aucQ}{0.561}
\newcommand{\aucQG}{0.588}
\newcommand{\aucQGHi}{0.607}
\newcommand{\aucQGLo}{0.570}
\newcommand{\aucQGPA}{0.593}
\newcommand{\aucQGPAHi}{0.611}
\newcommand{\aucQGPALo}{0.575}
\newcommand{\aucQGPASd}{0.004}
\newcommand{\aucQGPB}{0.610}
\newcommand{\aucQGPBHi}{0.628}
\newcommand{\aucQGPBLo}{0.592}
\newcommand{\aucQGPBSd}{0.004}
\newcommand{\aucQGPC}{0.619}
\newcommand{\aucQGPCHi}{0.637}
\newcommand{\aucQGPCLo}{0.601}
\newcommand{\aucQGPCSd}{0.003}
\newcommand{\aucQGSd}{0.003}
\newcommand{\aucQHi}{0.578}
\newcommand{\aucQLo}{0.543}
\newcommand{\aucQSd}{0.005}
\newcommand{\codeFlawed}{24}
\newcommand{\codeLowEnt}{6}
\newcommand{\codeN}{40}
\newcommand{\codeNotLoad}{5}
\newcommand{\codeQuest}{1}
\newcommand{\codeRepair}{4}
\newcommand{\dIncA}{+0.005}
\newcommand{\dIncAHi}{+0.014}
\newcommand{\dIncALo}{-0.004}
\newcommand{\dIncAP}{0.296}
\newcommand{\dIncB}{+0.022}
\newcommand{\dIncBHi}{+0.031}
\newcommand{\dIncBLo}{+0.011}
\newcommand{\dIncBP}{0.001}
\newcommand{\dIncC}{+0.031}
\newcommand{\dIncCHi}{+0.041}
\newcommand{\dIncCLo}{+0.020}
\newcommand{\dIncCP}{0.001}
\newcommand{\dPAvsQG}{-0.025}
\newcommand{\dPAvsQGHi}{-0.004}
\newcommand{\dPAvsQGLo}{-0.046}
\newcommand{\dPAvsQGP}{0.012}
\newcommand{\fatAll}{0.765}
\newcommand{\fatAllHi}{0.782}
\newcommand{\fatAllLo}{0.748}
\newcommand{\fatGsm}{0.966}
\newcommand{\fatGsmHi}{0.990}
\newcommand{\fatGsmLo}{0.942}
\newcommand{\fatLast}{0.890}
\newcommand{\fatMath}{0.842}
\newcommand{\fatMathHi}{0.869}
\newcommand{\fatMathLo}{0.814}
\newcommand{\fatOly}{0.756}
\newcommand{\fatOlyHi}{0.788}
\newcommand{\fatOlyLo}{0.725}
\newcommand{\fatOmni}{0.659}
\newcommand{\fatOmniHi}{0.692}
\newcommand{\fatOmniLo}{0.624}
\newcommand{\fatOnsetMax}{0.766}
\newcommand{\fatOnsetMin}{0.749}
\newcommand{\fatPredErrStep}{0.670}
\newcommand{\fatPredMaxInc}{0.003}
\newcommand{\fatPredPos}{0.484}
\newcommand{\fatPredQG}{0.671}
\newcommand{\lenFixA}{+0.052}
\newcommand{\lenFixAP}{0.032}
\newcommand{\markFat}{15.5}
\newcommand{\markP}{0.09}
\newcommand{\markRec}{18.6}
\newcommand{\mwRhoP}{0.21}
\newcommand{\mwRhoR}{0.04}
\newcommand{\nErr}{2,221}
\newcommand{\nFatal}{1,700}
\newcommand{\nRecovered}{521}
\newcommand{\nonFatalOmniPct}{34.1}
\newcommand{\nonFatalPct}{23.5}
\newcommand{\twErrFC}{0.709}
\newcommand{\twErrFCHi}{0.726}
\newcommand{\twErrFCLo}{0.691}
\newcommand{\twErrFR}{0.376}
\newcommand{\twErrFRHi}{0.402}
\newcommand{\twErrFRLo}{0.350}
\newcommand{\twErrRC}{0.819}
\newcommand{\twErrRCHi}{0.840}
\newcommand{\twErrRCLo}{0.796}
\newcommand{\twOutFC}{0.664}
\newcommand{\twOutFCHi}{0.683}
\newcommand{\twOutFCLo}{0.645}
\newcommand{\twOutFR}{0.456}
\newcommand{\twOutFRHi}{0.484}
\newcommand{\twOutFRLo}{0.428}
\newcommand{\twOutQGFC}{0.672}
\newcommand{\twOutQGFCHi}{0.692}
\newcommand{\twOutQGFCLo}{0.653}
\newcommand{\twOutQGFR}{0.501}
\newcommand{\twOutQGFRHi}{0.526}
\newcommand{\twOutQGFRLo}{0.473}
\newcommand{\twOutQGRC}{0.679}
\newcommand{\twOutQGRCHi}{0.706}
\newcommand{\twOutQGRCLo}{0.650}
\newcommand{\twOutRC}{0.703}
\newcommand{\twOutRCHi}{0.729}
\newcommand{\twOutRCLo}{0.677}
\newcommand{\wpG}{0.548}
\newcommand{\wpGHi}{0.591}
\newcommand{\wpGLo}{0.504}
\newcommand{\wpPA}{0.485}
\newcommand{\wpPAHi}{0.536}
\newcommand{\wpPALo}{0.440}
\newcommand{\wpPC}{0.539}
\newcommand{\wpPCHi}{0.589}
\newcommand{\wpPCLo}{0.487}
\newcommand{\wpPairs}{466}
\newcommand{\wpProblems}{326}
\newcommand{\wpQGPC}{0.586}
\newcommand{\wpQGPCHi}{0.634}
\newcommand{\wpQGPCLo}{0.535}

\definecolor{acc}{HTML}{2a78d6}
\definecolor{warn}{HTML}{eb6834}
\definecolor{ok}{HTML}{1baf7a}
\titleformat{\section}{\large\bfseries\color{acc}}{\thesection}{0.5em}{}
\titlespacing*{\section}{0pt}{7pt}{3pt}
\setlength{\parskip}{3pt}\setlength{\parindent}{0pt}
\newtcolorbox{idea}[1][]{colback=acc!5,colframe=acc,fonttitle=\bfseries\small,title={#1},left=4pt,right=4pt,top=2pt,bottom=2pt,boxsep=2pt}
\newtcolorbox{dhyan}[1][]{colback=warn!6,colframe=warn,fonttitle=\bfseries\small,title={#1},left=4pt,right=4pt,top=2pt,bottom=2pt,boxsep=2pt}
\renewcommand{\arraystretch}{1.12}

\begin{document}
\begin{center}
{\LARGE\bfseries Errors Are Not Choke Points: poori research ek kahani mein}\\[3pt]
{\small Harsh Dubey, Harshit Saini, Prashant Jain, Himanshu Kumar, Gaurav Kaushal, Rohit Chaudhary (supervisor) \;·\; KIET, CSE (AI\&ML)}\\[1pt]
{\small Simple Hinglish guide. Ise padhne ke baad paper (\texttt{paper/main.pdf}) ek baar padho, sab samajh aa jayega.}
\end{center}

\section{Shuruaat: problem kya tha?}
LLM (jaise Qwen, Llama) maths ka sawaal step-by-step solve karta hai. Isko \textbf{chain of thought (CoT)}
ya \textbf{trace} kehte hain. Kabhi beech mein ek step galat ho jaata hai aur final answer galat aa jaata hai.
Poori field ka common idea yeh hai:
\begin{enumerate}
\item Pehla galat step dhundho (isko \textbf{first error} kehte hain; human annotators isko mark karte hain).
\item Wahan tak wapas jao aur wahan se dobara generate karo (\textbf{backtracking / repair}).
\end{enumerate}
Process Reward Models (PRMs) isi kaam ke liye train hote hain, aur ProcessBench jaise benchmarks isi ko grade
karte hain: ``kya tumne pehla galat step pakda?'' Hamara original aim tha \textbf{early failure detection}:
kya hum trace ke shuru mein hi bata sakte hain ki yeh fail hoga, taaki jaldi intervene kar sakein?

\begin{center}\small
\begin{tabularx}{\textwidth}{lX}
\toprule
Phase & Kahani ek nazar mein \\
\midrule
1. Early detection & detector bana, par woh sirf sawaal ki difficulty padh raha tha (negative result) \\
2. Naya sawaal, $v_k$ & failure kahan decide hota hai, model ke apne continuations se mapa \\
3. Main result & error $\ne$ choke point; repair, restart ko same compute par nahi harata; theory se proof \\
4. ChokeBench & naya benchmark; check kiya ki yeh asli recovery predict karta hai \\
5. RAG (secondary) & hallucination par same idea; result grader par depend karta hai, human rating baaki \\
6. Review aur refocus & do reviews ke baad claims narrow kiye, kamzor hisse imaandari se likhe \\
7. Run D & reasoning models (Qwen3, R1) par exact on-policy test, Kaggle par chal raha hai \\
\bottomrule
\end{tabularx}
\end{center}

\section{Phase 1: early detection, aur pehla jhatka}
ProcessBench (3,400 solutions, 12 models, har solution mein human-marked first error) par humne detectors
banaye. Number achhe lag rahe the, lekin phir humne ek simple \textbf{baseline} lagaya: sirf sawaal ka text
aur kaunsa model hai, trace ka ek bhi step dekhe bina.

\begin{center}\small
\begin{tabular}{lcl}
\toprule
Detector ko kya diya & AUROC & Matlab \\
\midrule
Sirf sawaal + model (koi step nahi) & \aucQG & sawaal kitna mushkil hai, bas itna pata \\
+ trace ke pehle 25\% steps & \dIncA\ (n.s.) & early steps se lagbhag kuch nahi milta \\
+ poora trace & \dIncC & thoda signal, par tab tak bahut der ho chuki \\
\bottomrule
\end{tabular}
\end{center}
\begin{dhyan}[Seekh (negative result)]
Early detector asal mein \emph{``sawaal mushkil hai''} padh raha tha, \emph{``yeh attempt fail hoga''} nahi.
Isko \textbf{difficulty confound} kehte hain. Yahi hamara pehla pivot tha: naya sawaal poocho.
\end{dhyan}

\section{Phase 2: naya sawaal, aur ek hi measurement $v_k$}
Naya sawaal: \textbf{failure asal mein kahan ``decide'' hota hai, aur kya wahan se repair karna restart se behtar hai?}
Iske liye ek quantity:
\[ v_k = P(\text{sahi answer} \mid \text{sawaal} + \text{pehle } k \text{ steps}) \]
Isko estimate karne ka tareeka: trace ko step $k$ par kaato, \emph{usi model} se wahan se 8--16 baar aage
likhwao (\textbf{continuations / rollouts}), aur gino kitni baar answer sahi aaya. Total \TotalRollouts\
continuations chalaye (Kaggle ke free T4 GPUs par).

\begin{center}\small
\begin{tabularx}{\textwidth}{lX}
\toprule
Term & Simple matlab \\
\midrule
$v_0$ & bilkul shuru se solve karne ka chance. Yahan se dobara likhna = \textbf{restart}. \\
Backtrack to step $j$ & step $j$ se just pehle tak rakho, wahan se dobara likho (yaani $v_{j-1}$ wala point). \\
\textbf{Choke point} (decisive step) & woh step jahan $v$ sabse zyada girta hai. Yahan failure asal mein decide hota hai. \\
\textbf{First error} & woh step jahan human kehta hai ``yahan galti hui''. \\
Capability-limited & $v_0 \le 1/8$: model ke liye sawaal shuru se hi lagbhag namumkin. \\
\bottomrule
\end{tabularx}
\end{center}
\textbf{Example:} $v = [0.6, 0.6, 0.1, 0.0]$. Step 2 par $v$ 0.6 se 0.1 gira, toh choke point step 2 hai.
Agar human ne step 3 ko galat mark kiya, toh \emph{error} aur \emph{choke point} alag hain.

\section{Phase 3: kya mila (main results, maths)}
\begin{center}\small
\begin{tabularx}{\textwidth}{>{\raggedright\arraybackslash}p{5.3cm}>{\raggedright\arraybackslash}p{4.6cm}X}
\toprule
Sawaal & Jawab (number) & Matlab \\
\midrule
Kitne failures step 1 se pehle hi decide? & \QwenSmallAllCapPct--\QwenMathCapPct\% (4 models) & repair karne ko koi step hi nahi \\
Baaki mein human error = choke point? & \QwenTwoAllDec--\QwenSmallAllDec\ (chance 0.16--0.20) & chance se behtar, par aadha hi sahi \\
Naye model (Llama-3.1-8B) par confirm? & haan, \LlamaDec\ & preregistered test, pass \\
PRM kya pakadta hai? (ChokeBench) & human error \cbPrmDropAccHPct\%, choke point sirf \cbPrmDropAccDPct\% & verifier label pakadta hai, asli point nahi \\
Choke point ke liye train kiya model? & \cbLearnedAccDPct\% & better nahi hua \\
Restart kitna bachata hai? (4 attempts) & \recMinFour--\recMaxFour\% failures solve & simple restart bahut strong \\
PRM backtracking vs restart (same compute, Llama) & \LlamaPolPrmDropFourDiff\ (confirmed) & repair, restart se \emph{haara} \\
Free perfect oracle vs restart (4$\times$ budget) & sirf \cbOracleGainB & jaadu ki chhadi bhi bas thoda jeetti \\
\bottomrule
\end{tabularx}
\end{center}
\textbf{Matched compute} ka matlab: har method ko barabar GPU budget (tokens, PRM ki cost sab gini gayi).
Tabhi comparison fair hai.

\begin{idea}[Kyun? (theory, Proposition 1: supermartingale)]
Ek failed trace ka koi bhi point, jo aage dekhe bina chuna gaya ho, average mein shuru se behtar nahi hota:
$E[v_\tau \mid \text{fail}] \le v_0$. Simple bhasha: \emph{jo trace fail hua, uska aadha hissa ek fresh start
se zyada keemti nahi hota}. Isliye choke point dhundhne se tokens bachte hain, success nahi badhti.
(Proposition 2: iska fayda bas chhote attempts hai, budget ka lagbhag 1.25--1.4$\times$.)
\end{idea}

\section{ChokeBench: hamara naya benchmark}
ProcessBench poochta hai ``pehla galat step pakda?''. \textbf{ChokeBench v1} (\cbN\ failed traces, 4 models)
\emph{dono} poochta hai: human error pakda, \emph{aur} choke point pakda? Aur woh step resampling ke liye kitna keemti hai?
Sabse important check: \textbf{kya ChokeBench score asli recovery predict karta hai?} (circular na ho, isliye samples ke
aadhe hisse se choke point label kiya, baaki aadhe se gain mapa.)
\begin{center}\small
\begin{tabular}{lc}
\toprule
Check (exploratory) & Result \\
\midrule
Named step = choke point: recovery gain mein badhat & \valBothD \\
Named step = human error, choke point ko control karne ke baad & \valBothE\ (n.s.) \\
Naye (held-out) model par ranking, Spearman $\rho$: ChokeBench vs ProcessBench-style & \valRhoD\ vs \valRhoE \\
\bottomrule
\end{tabular}
\end{center}
Matlab: human label tabhi kaam aata hai jab woh choke point se match kare. ChokeBench asli fayde ko behtar predict karta hai.

\section{Phase 4 (secondary): hallucination wala RAG study}
RAG mein model ko ek source diya jaata hai aur answer source se hi aana chahiye. Source mein jo nahi hai
woh likh diya toh \textbf{hallucination}. Data: RAGTruth (human labels), model Llama-2-7B. Detector $D$ =
LettuceDetect (har sentence ko score karta hai).
\begin{center}\small
\begin{tabularx}{\textwidth}{>{\raggedright\arraybackslash}p{5.6cm}cX}
\toprule
Method & Cost (base = 1$\times$) & Judged hallucination change vs base (E1 judge) \\
\midrule
DRA: detect, restart, abstain & \costDra$\times$ & \vbDraEOne\ (exploratory) \\
CARE (choke-point-aware controller) & \costCare$\times$ & \rrHZero\ ($p=\rrHZeroP$, preregistered) \\
CARE vs plain restart & -- & \rrHOne\ (koi fark nahi) \\
Chain-of-Verification & \costCove$\times$ & mehenga, aur behtar nahi \\
\bottomrule
\end{tabularx}
\end{center}
\textbf{Proposition 3 (law):} restart loop jo answer deta hai, uski hallucination rate = detector ka
$P(\text{hallucinated}\mid\text{pass})$, retries kitne bhi ho. Retries sirf coverage badhaate hain. Human
labels se prediction: \lawH\% $\to$ \lawPi\%. Check: 4$\times$ retries par coverage \invCov\ points, hallucination
lagbhag same (\invETwo\ se \invEOne). Prediction sahi nikli.
\begin{dhyan}[Imaandari: yeh part kamzor kyun hai]
(1) Judge E1 humans se kam agree karta hai ($\kappa=\evEOneKappa$; 81\% drafts flag karta hai jab humans 50\%).
(2) Grader badlo toh result badal jaata hai: detector jaisa grader \robBaseETwoDiff, independent judge \robBaseEOneDiff\ (n.s.).
Gate khud ko grade kar raha tha. (3) Outputs ko abhi tak humans ne rate nahi kiya. Isliye RAG
\emph{secondary} hai, aur hum ``hallucination solve'' ka claim \textbf{nahi} karte.
\end{dhyan}

\section{Imaandari ke kadam (reviewers yahi dekhte hain)}
\begin{center}\small
\begin{tabularx}{\textwidth}{lX}
\toprule
Kya kiya & Kyun \\
\midrule
\textbf{Preregistration} & hypotheses data se pehle git mein commit kiye, taaki baad mein result chun na sakein \\
\textsc{pre} / \textsc{exp} / \textsc{conf} tags & har claim ke saath likha: pehle se plan tha, baad mein dekha, ya naye data par confirm hua \\
Sample splitting & ek half se chuno, doosre se score karo; warna restart nakli tareeke se achha dikhta \\
Bootstrap CIs, Holm correction & har number ke saath uncertainty; kai tests ki galti control \\
Mistakes withdraw kiye & ``dual judge'' aur ``1,500-item'' table galat the, hata diye, corrections log mein likha \\
Blinded human kits & base vs DRA ke 80 pairs, 2 raters; key ka hash pehle commit, key repo se bahar \\
External review & do reviews ke baad focus narrow kiya: main claim maths wala, RAG secondary \\
\bottomrule
\end{tabularx}
\end{center}

\section{Run D: naya bada experiment (abhi Kaggle par chal raha hai)}
\textbf{Kyun?} Pichhle traces ProcessBench ke the, kisi aur ne likhe the, isliye sirf ``near on-policy'' the.
Proposition 1 exactly on-policy traces par hold karti hai. Aur aaj ki field ka sawaal \emph{reasoning models}
ka hai (jo hazaaron tokens sochte hain, \texttt{<think>} ke andar khud backtrack karte hain). Run D yahi test karta hai.
\begin{center}\small
\begin{tabularx}{\textwidth}{lX}
\toprule
Setting & Value \\
\midrule
Models & Qwen3-1.7B (thinking mode), DeepSeek-R1-Distill-Qwen-1.5B \\
Data & MATH-500, problems 21--500 (480; pehle 20 sirf pilot ke liye, analysis se bahar) \\
Stage 1 & har problem ka ek trace, max 6,144 tokens; math-verify se grade \\
Stage 2 & har failed trace ko 9 jagah kaato: 0, 256, 512, 1024, 1536, 2048, 3072, 4096, 5120 tokens
(paragraph break par snap, taaki cut sirf past par depend kare); har cut se $m=8$ continuations \\
Hardware & Kaggle, 2$\times$ T4 GPU, vLLM; har model ka $\sim$6.5 ghante ka session \\
\bottomrule
\end{tabularx}
\end{center}
\begin{center}\small
\begin{tabularx}{\textwidth}{lX}
\toprule
Hypothesis (preregistered) & Test \\
\midrule
D1: Proposition 1 exact on-policy par & har cut $t$ par $v_t - v_0 \le 0$? Agar kisi cut par significantly positive (Holm $p<0.05$), toh theory refute \\
D2: koi cut restart ko 5 points se nahi harata & same token budget par best cut (half A se chuna, half B par score) minus restart; 4$\times$ par upper CI $< +0.05$ \\
D3: descriptive & kitne failures shuru se decided ($v_0 \le 1/8$), aur $v$ sabse zyada kahan girta hai \\
D4: sensitivity & D1, D2 bina un traces ke jo token limit par kat gaye \\
\bottomrule
\end{tabularx}
\end{center}
\textbf{Ab tak:} pilot (8 problems) mein pipeline sahi chali; 6/8 sahi, dono failures token limit par kate.
Speed kam thi ($\sim$130 tokens/s per GPU), isliye fast CUDA-graph mode lagaya, $m=8$ rakha, aur problems dataset
order mein time tak process honge. Yeh sab \textbf{deviations data aane se pehle} preregistration doc mein likhe.
Qwen3 ka run chal raha hai (raat ~2:40 baje khatam), uske baad R1 shuru hoga (Kaggle ek waqt mein ek hi slot de raha hai).
Result jo bhi aaye (theory sahi ya galat), paper mein waise hi jayega. Agar D2 mein koi cut restart ko
harata hai, toh hum likhenge ki reasoning models conclusion badalte hain.

\section{Ant mein: ek line mein research}
\begin{idea}[Final takeaway]
\textbf{Errors are not choke points.} Jahan human galti dikhata hai, failure aksar wahan decide nahi hota;
bahut failures shuru se hi decided hote hain; verifiers label pakadte hain, asli point nahi; aur us point se
repair karna, same compute par, simple restart ko nahi harata. Isliye restart woh baseline hai jise har
localise-and-repair method ko harana chahiye. Hamara naya ChokeBench yeh test karne ka tool hai.
\end{idea}
\small\textbf{Abhi baaki:} (1) Run D ke results; (2) human ratings (base vs DRA, 2 log, ~2--3 ghante each);
(3) doosre researchers ka replication. \textbf{Paper padhne ka order:} Abstract $\to$ Intro ke 5 contributions $\to$
Sec.~5--7 (counterfactual, policies, ChokeBench) $\to$ Sec.~8--9 (RAG) $\to$ Limitations $\to$ Appendix (proofs).
\end{document}
